\SourcePage{36}{39}
\section{The Wave Function and Measurements}

Let us return to the measurement process, whose properties were discussed
qualitatively in \secref{1}, and show how those properties are connected with the
mathematical apparatus of quantum mechanics.

Consider a system made up of two parts: a classical apparatus and an electron
(treated as a quantum object). Measurement consists in bringing these two
parts into interaction. The apparatus consequently passes from its initial
state into some other state, and from this change we draw conclusions about
the state of the electron. States of the apparatus are distinguished by the
values of a physical quantity characterizing it (or%
\SourcePage{37}{40}
quantities)---the apparatus ``reading.'' Denote this quantity provisionally by
$g$, and its eigenvalues by $g_n$. Because the apparatus is classical, these
values generally range over a continuum; solely to simplify the notation in
the formulas below, however, we shall suppose that the spectrum is discrete.
The states of the apparatus are described by quasiclassical wave functions,
which we denote by $\Phi_n(\xi)$; the index $n$ corresponds to the reading
$g_n$, and $\xi$ stands collectively for its coordinates.
That the apparatus is classical manifests itself in the circumstance that at
each instant one can assert with certainty that it occupies one of the known states
$\Phi_n$, with some definite value of $g$. Such a statement would of course
be false for a quantum system.

Let $\Phi_0(\xi)$ be the wave function of the apparatus in its initial state
(before measurement), and let $\Psi(q)$ be an arbitrary normalized initial
wave function of the electron ($q$ denotes its coordinates). These functions
describe apparatus and electron independently, so the initial wave function
of the complete system is the product
\begin{equation}
  \Psi(q)\Phi_0(\xi).
  \tag{7.1}
\end{equation}
The apparatus and electron then interact. In principle, the equations of
quantum mechanics allow the time variation of the system's wave function to
be followed. After the measurement it will plainly no longer be a product of
functions of $\xi$ and $q$. Expanding it in the eigenfunctions $\Phi_n$ of the
apparatus (which form a complete system of functions), we obtain a sum of the
form
\begin{equation}
  \sum_n A_n(q)\Phi_n(\xi),
  \tag{7.2}
\end{equation}
where the $A_n(q)$ are certain functions of $q$.

At this point the ``classical nature'' of the apparatus enters, together with
the double role of classical mechanics as both a limiting case and a basis of
quantum mechanics. As already noted, because the apparatus is classical, the
quantity $g$ (its ``reading'') has a definite value at every instant. It
follows that after measurement the actual state of apparatus plus electron is
described not by the whole sum (7.2), but by only the one term corresponding
to the apparatus reading $g_n$:
\begin{equation}
  A_n(q)\Phi_n(\xi).
  \tag{7.3}
\end{equation}
\SourcePage{38}{41}
Consequently, $A_n(q)$ is proportional to the electron's wave function after
measurement. It is not itself that wave function, as is already evident from
the fact that $A_n(q)$ is not normalized. It contains both information about
the properties of the state produced and the probability, determined by the
system's initial state, that the apparatus will give its $n$th reading.

By the linearity of the equations of quantum mechanics, the relation between
$A_n(q)$ and the initial electron wave function $\Psi(q)$ is in general given
by a linear integral operator,
\begin{equation}
  A_n(q)=\int K_n(q,q')\Psi(q')\,dq',
  \tag{7.4}
\end{equation}
whose kernel $K_n(q,q')$ characterizes the measurement process in question.

We suppose that the measurement under consideration yields a complete
description of the electron's state. In other words (see \secref{1}), in the state
produced, the probabilities for all quantities must be independent of the
electron's preceding state (before measurement). Mathematically, this means
that the form of the functions $A_n(q)$ must be fixed by the measurement
process itself and must not depend on the initial electron wave function
$\Psi(q)$.

The $A_n$ must therefore have the form
\begin{equation}
  A_n(q)=a_n\varphi_n(q),
  \tag{7.5}
\end{equation}
where the $\varphi_n$ are certain functions that we shall take to be
normalized, while only the constants $a_n$ depend on the initial state
$\Psi(q)$. In the integral relation (7.4), this corresponds to a kernel
$K_n(q,q')$ which factors into functions of $q$ and $q'$ separately:
\begin{equation}
  K_n(q,q')=\varphi_n(q)\Psi_n^*(q').
  \tag{7.6}
\end{equation}
The linear dependence of the constants $a_n$ on $\Psi(q)$ is then expressed
by formulas of the form
\begin{equation}
  a_n=\int\Psi(q)\Psi_n^*(q)\,dq,
  \tag{7.7}
\end{equation}
where the $\Psi_n(q)$ are certain functions fixed by the measurement process.

The functions $\varphi_n(q)$ are the normalized wave functions of the
electron after measurement. We thus see how the theory's mathematical
formalism represents the possibility of producing, by a measurement, an
electron state described by a definite wave function.
\SourcePage{39}{42}
When a measurement is made on an electron having a specified wave function
$\Psi(q)$, the constants $a_n$ have a simple physical meaning: by the general
rules, $|a_n|^2$ is the probability that the measurement gives the $n$th
result. The probabilities of all results sum to unity:
\begin{equation}
  \sum_n|a_n|^2=1.
  \tag{7.8}
\end{equation}
For an arbitrary normalized function $\Psi(q)$, the validity of (7.7) and
(7.8) is equivalent (compare \secref{3}) to the statement that $\Psi(q)$ can be
expanded in the functions $\Psi_n(q)$. Hence the functions $\Psi_n(q)$ form
a complete normalized and mutually orthogonal set.

If the electron's initial wave function coincides with one of the functions
$\Psi_n(q)$, the corresponding constant $a_n$ is plainly unity and all the
others vanish. Put differently, if the measurement is performed on an electron
in the state $\Psi_n(q)$, it will with certainty give a definite ($n$th) result.

Taken together, these properties of the functions $\Psi_n(q)$ demonstrate
that they are eigenfunctions of a certain physical quantity that characterizes the electron (call it
$f$), and the measurement being considered may be called a measurement of
that quantity.

It is essential that, in general, the functions $\Psi_n(q)$ do not coincide
with the functions $\varphi_n(q)$ (indeed, the latter need not even be
mutually orthogonal or form an eigenfunction system of any operator). This
fact expresses first of all the non-repeatability of measurement results in
quantum mechanics. If the electron was in the state $\Psi_n(q)$, a measurement
of $f$ made upon it will reveal the value $f_n$ with certainty. After the
measurement, however, the electron is in a different state $\varphi_n(q)$, in
which $f$ need no longer have any definite value. Thus, if a second
measurement were made immediately after the first, the value found for $f$
would not coincide with the result of the first
measurement\footnote{There is, however, an important exception to the
non-repeatability of measurements: coordinate is the one quantity whose
measurement is repeatable. Two measurements of an electron's coordinate made
at a sufficiently short interval must give nearby values; otherwise the
electron would have infinite velocity. The mathematical origin of this
exception lies in the commutativity of the coordinate with the
electron--apparatus interaction energy operator, which (in the nonrelativistic
theory) depends solely on
coordinates.}. To predict (in the sense of%
\SourcePage{40}{43}
calculating the probability of) the result of a repeated measurement when the
first result is known, one takes from the first measurement the wave function
$\varphi_n(q)$ of the state it produced, and from the second the wave function
$\Psi_m(q)$ of the state whose probability is sought. More explicitly, the
equations of quantum mechanics determine a wave function
$\varphi_n(q,t)$ which equals $\varphi_n(q)$ at the time of the first
measurement. If the second measurement is made at time $t$, the probability
of its $m$th result is the squared modulus of the integral
$\int\varphi_n(q,t)\Psi_m^*(q)\,dq$.

We see that measurement in quantum mechanics is ``two-faced'': its roles with
respect to past and future do not coincide. Toward the past it ``verifies''
the probabilities of the various possible results predicted from the state
created by the preceding measurement. Toward the future it creates a new
state (see also \secref{44}). A profound irreversibility is therefore inherent in
the measurement process itself.

This irreversibility has important consequences of principle. As will be seen
later (see the end of \secref{18}), the fundamental equations of quantum mechanics
by themselves are symmetric under reversal of the sign of time; in this
respect quantum mechanics does not differ from classical mechanics. The
irreversibility of measurement, however, makes the two directions of time
physically inequivalent in quantum phenomena, and thus produces a distinction
between future and past.
